\section{引言}
\subsection{研究背景}
在全球“双碳”战略背景下，科学识别碳排放驱动因素并预测达峰路径是核心议题。碳排放受经济规模、产业结构、能源与技术等多因素交织影响。在“双碳”目标下，理清这些因素的联动机制并预测达峰路径，是当前亟待解决的现实问题。

从研究尺度看，省级面板分析能够呈现区域差异，但在小样本宏观时序数据与统计口径一致性约束下，往往面临变量定义不统一、口径协调成本高、结果解释分散等问题。与之相比，国家级时间序列分析可以在较低的数据整合成本下形成统一口径，更适合回答总量控制、达峰时点和政策路径等宏观问题。在面向公共政策评估的研究语境中，国家级框架更便于形成“数据清洗—机制识别—预测比较—情景模拟—政策解释”的完整链条，并在有限数据可得性条件下保持结论可复核性。

\subsection{研究问题与研究意义}
围绕国家层面碳达峰路径识别，本文聚焦三个递进问题：第一，人口、富裕度、能源强度、能源结构、产业结构、技术进步和治理投入等变量对碳排放的方向与强度如何；第二，在国家级小样本时间序列条件下，不同预测模型的外推能力是否存在稳定差异；第三，在不同政策强度下，未来碳排放峰值水平与达峰路径将出现何种分化。上述问题分别对应“机制解释”“预测适配”“政策模拟”三个层次，构成本文的主线。

从现实层面看，国家级碳达峰预测有助于评估减排路径的可行性，并为能源结构优化、产业升级、技术创新与环境治理政策提供量化证据。相比只给出总量趋势，机制识别能够进一步回答“为什么会出现该趋势”，从而提高政策建议的针对性。对当前“双碳”政策实践而言，准确识别煤炭依赖、工业结构、技术扩散和治理投入在排放变化中的相对作用，是优化政策组合的重要前提。

从方法层面看，将解释性模型与预测性模型结合，可以避免“只解释不预测”或“只预测不解释”的单一化研究倾向。扩展STIRPAT和EKC检验强调变量作用机制，中介分析强调结构传导路径，滚动一步预测强调真实外推能力，多情景模拟强调政策含义外化。多方法协同不仅提升研究严谨性，也增强了结果在公共政策评估语境中的可检验性和可沟通性。

\subsection{研究思路、创新点与边界}
本文研究思路可概括为“数据整合—机制识别—预测比较—情景模拟—政策解释”五个环节。首先，基于CEADs省级表观碳排放清单按年汇总构建国家级碳排放序列，并与国家统计局宏观指标进行口径匹配；其次，通过EKC检验和扩展STIRPAT模型识别核心驱动因素及其弹性方向；再次，利用中介框架讨论能源结构、技术进步和产业结构之间的作用链；随后，在滚动一步预测框架下比较多种模型的外推表现；最后，设定基准、低碳和深度脱碳三类情景并开展峰值路径模拟，形成政策含义。

本文创新主要体现在四个方面。第一，构建了“机制解释—预测比较—情景模拟”一体化的分析框架，实现了从现象解释到政策推演的连贯分析。第二，在国家级时序框架内，将能源结构、产业结构、技术进步与治理投入纳入统一分析体系，提升变量覆盖的现实对应性。第三，采用滚动一步预测评估模型而非仅比较样本内拟合，强化了对外推能力的识别。第四，结合滚动预测误差与情景偏离度，构建了政策效果的动态监测思路，为后续减排参数的年度调整提供量化参考。

需要说明的是，本文结论适用于国家总量层面，不直接推断区域异质性和行业细分机制；2023—2024年数据用于预测衔接与情景延展，不参与机制回归估计。该边界设置有助于保证统计推断的口径一致性，避免将预测值误作观测事实。
